\documentclass[11pt]{article}
\usepackage[a4paper,margin=25mm]{geometry}
\usepackage{amsmath,amssymb,amsthm}
\usepackage[hidelinks]{hyperref}
\emergencystretch=2em
\newtheorem{theorem}{Theorem}
\title{Critical sparse random graphs violate a forest axiom\newline
\large Disproof of Conjecture 00000005029}
\author{gaochengzhecpu}\date{}
\begin{document}\maketitle
\begin{abstract}
For every $n$ divisible by three, the genuine random
graph $G(n,1/n)$ contains a triangle with probability
at least $1/50$. Thus a single first-order forest
axiom, the absence of triangles, has probability at
most $49/50$ on an unbounded sequence of sizes. It
cannot belong to the asserted almost-sure forest
limiting theory. The proof uses actual independent
edge probabilities and an elementary ordered-pair
Bonferroni bound; no Poisson limit or random-graph
asymptotic theorem is needed.
\end{abstract}

\section*{The sentence that must hold for a forest theory}
In the language of simple graphs, consider
\[
 \varphi=\forall x\,\forall y\,\forall z\,
       \neg\bigl(E(x,y)\wedge E(y,z)\wedge E(z,x)\bigr).
\]
The three adjacencies automatically force the vertices
to be distinct, since the graph has no loops.
Every forest satisfies this first-order sentence.
If the forest axioms formed the limiting theory in
the asserted sense, the probability of $\varphi$
would tend to one.

We bound this single sentence's probability directly.
This avoids exchanging the limit with an infinite
conjunction of all cycle exclusions. A bound for the
actual forest event will also follow, but it is not
needed for the sentence-level contradiction.

\section*{The actual independent-edge model}
Let $h\geq1$, put $n=3h$, and use vertices
$\{0,1,2\}\times\{0,\ldots,h-1\}$. Each unordered
pair of distinct vertices is independently present
with probability $p=1/n$. This is exactly $G(n,1/n)$,
the source's $c=1$ model. No conditioning is imposed.

For an explicit finite construction, choose for every
unordered pair an independent uniform label in
$\{0,\ldots,n-1\}$ and retain the edge precisely
when the label is zero. An arbitrary finite prescribed
pattern of present and absent edges then has probability
$p^{a}(1-p)^{b}$, where $a,b$ count its present and
absent entries. The Lean project proves this complete
pattern formula for the actual mapped graph PMF.

\newpage
\section*{Triangles in three vertex blocks}
For each triple $t=(a,b,c)$ with one vertex in each
of the three blocks, let $E_t$ mean that its three
actual edges are present. There are $h^3$ triples.
The edges are distinct, so independence gives
\[
               \Pr(E_t)=p^3,
           \qquad \sum_t\Pr(E_t)=h^3(3h)^{-3}=\frac1{27}.
\]
Two distinct selected triangles can share at most
one edge. If they share an edge, their union has
five distinct edges and their joint probability is
$p^5$. Otherwise the union has six edges and the
joint probability is $p^6$, even when a vertex is shared.

There are at most $h^6$ ordered pairs of triangles
of the latter kind. For shared edges, overcount
separately by each of the three block-pair types.
Sharing an edge fixes the same two corresponding
coordinates in the two triples. Choosing those two
coordinates and the two remaining coordinates
leaves $h^4$ possibilities. Thus there are at most
$3h^4$ ordered pairs sharing an edge. Harmless
overcounts, including identical triples in this
counting bound, only increase the upper estimate.
It follows that
\[
 \sum_{t\ne s}\Pr(E_t\cap E_s)
       \leq h^6p^6+3h^4p^5
       =\frac1{729}+\frac1{81h}.
\]
All sums here count ordered pairs, not unordered pairs.

\section*{An ordered-pair Bonferroni lower bound}
For any finite family of events, the deliberately
weakened second Bonferroni inequality is
\[
 \Pr\Bigl(\bigcup_t E_t\Bigr)
     \geq\sum_t\Pr(E_t)
                  -\sum_{t\ne s}\Pr(E_t\cap E_s).
\]
To prove it pointwise, let $k$ be the number of events
that occur. If $k=0$, all terms are zero. If $k\geq1$,
then $k\leq1+k(k-1)$, so the sum of indicators is at
most the union indicator plus the ordered distinct-pair
indicator sum. Taking expectations gives the inequality.
This intentionally subtracts both orientations; no
factor $1/2$ is missing from the estimate being used.

Combining the two bounds and $h\geq1$ gives
\[
 \Pr\Bigl(\bigcup_t E_t\Bigr)
   \geq\frac1{27}-\frac1{729}-\frac1{81h}
   \geq\frac{17}{729}>\frac1{50}.
\]
Every selected event witnesses a genuine triangle in
the actual graph, so the same lower bound applies
to the event that the graph has any triangle.

\newpage
\section*{The asymptotic contradiction}
For every $h\geq1$, therefore,
\[
 \Pr\bigl(G(3h,1/(3h))\models\varphi\bigr)
                    \leq\frac{49}{50}.
\]
The sizes $3h$ tend to infinity. If the full-size
sentence probability tended to one, its restriction
to this unbounded subsequence would also tend to one,
contradicting the fixed upper bound.

\begin{theorem}
At critical sparsity $c=1$, the forest axioms cannot
be the asserted almost-sure limiting theory of
$G(n,c/n)$.
\end{theorem}
\begin{proof}
The fixed first-order forest sentence $\varphi$ has
probability at most $49/50$ along every size $n=3h$.
It therefore fails the necessary probability-one
limit for membership in that limiting theory.
\end{proof}

There is also a direct bound for the actual forest
event. Three pairwise adjacent vertices produce a
genuine closed walk of length three whose intermediate
vertices are distinct. It is a cycle, so the graph
is not acyclic. Consequently
\[
 \Pr\bigl(G(3h,1/(3h))\text{ is a forest}\bigr)
       \leq\Pr\bigl(G(3h,1/(3h))\models\varphi\bigr)
       \leq\frac{49}{50}.
\]
The Lean project proves this bridge using an actual
Mathlib cyclic walk and its actual acyclicity predicate.
It separately proves that neither probability tends
to one and that the vertex cardinality tends to infinity.

The source is a conjunction. The false critical
forest-theory assertion suffices to refute it;
the separate supercritical giant-component and
quantifier-depth assertions are outside this submission.
No claim that the full graph sequence lacks a
probabilistic first-order convergence law is made.

\section*{Independent exact checks}
The accompanying \texttt{verify.py} enumerates all
eight simple graphs on three vertices and all
32,768 simple graphs on six vertices, weighting
every edge subset by its exact $G(n,1/n)$ mass.
It independently checks the triangle-event bound.
For block sizes $1$ through $8$, it also enumerates
all ordered selected-triangle pairs and verifies
their five/six-edge union cases and exact overlap
counts, using only integer and rational arithmetic.
These are supplementary consistency checks; Lean
proves the result for every positive block size.

\newpage
\section*{Formalization and reproducibility}
The formal vertex set is
$\operatorname{Fin}(3)\times\operatorname{Fin}(m+1)$,
with its proved cardinality $3(m+1)$. Uniform assignments
of labels to unordered pairs are mapped to genuine
Mathlib simple graphs. Auxiliary diagonal coordinates
do not affect this map: it explicitly forbids loops.
Actual event probabilities are taken from the graph
PMF, with a proved equality to normalized counting
on its original sample space.

The complete present/absent finite-edge pattern law
is derived by counting allowed coordinate values.
Each present coordinate has one allowed label; each
absent coordinate has $n-1$. Normalizing their product
proves the independent Bernoulli law on every actual
nonloop edge family, rather than only on selected
three-edge marginals.

The selected triangle event uses its actual three
unordered edges, proved distinct. Unequal triangle
indices have union cardinality five or six by symbolic
coordinate cases. The overlap counts are actual finite
double sums, and the ordered-pair Bonferroni inequality
is proved from genuine event indicators. These yield
the actual graph triangle lower bound, the literal
three-adjacency sentence bound and the forest bound.

Lean formalizes the graph property that is exactly
the semantics of the displayed sentence, not separate
first-order syntax. The mathematical sentence-level
argument is given above; it uses only that single
necessary forest axiom, not an interchange of an
infinite conjunction with probability limits.

Lean 4.19.0 and Mathlib revision\\
\texttt{c44e0c8ee63ca166450922a373c7409c5d26b00b}
are pinned. From \texttt{lean/}, run
\begin{verbatim}
lake exe cache get
lake build
lake env lean -DwarningAsError=true Main.lean
\end{verbatim}
Fresh-build logs, source and hashes are included.
Printed dependencies contain only standard foundational
axioms. Compile with \texttt{tectonic main.tex}.
The author agent and parent reviewed the work internally;
no external independent review is claimed. The parent
developed the exact finite verifier separately.

\begin{thebibliography}{9}
\bibitem{source} The Last Math Competition,
Conjecture 00000005029. Exact bilingual source included
as \texttt{SOURCE.md}. The finite probability bound
and its asymptotic implication are proved directly
above and in the accompanying Lean project.
\end{thebibliography}
\end{document}
